\scanpage{251}
\mspage{126}

\uline{21 Connection with Teichmüller spaces.}

First a supplement connected with the diagram \inserted{of inclusions} (31) of number 19, in which \(A_{g,(\nu,\nu+1)}\) deserves the simpler notation \(A_{g,\nu}^{\bullet}\) (\unsure{group} of automorphisms of \(U_{g,\nu}\) pointed at \(a_{g,\nu}\)). Consider \unsure{the inclusion|the filtration} of subgroups of \(A_g\)
\[
 A_{g,\nu}^{\circ}\supset B_{g,\nu}\supset A_{g,\nu+1}^{\circ}
 \tag{72}
\]
\deleted{\unsure{\(A_{g,\nu}^{\bullet}\)}}

\unsure{giving}, on dividing by \(A_{g,\nu+1}\) and in the anabelian case \(2g+\nu\geq 3\), the fibration

(73) \diagram{\begin{tikzcd}[ampersand replacement=\&]
\widetilde U_{g,\nu} \arrow[r] \& U_{g,\nu} \\
A_{g,\nu}^{\circ}/A_{g,\nu+1}^{\circ} \arrow[u, phantom, "\unsure{\simeq}" description] \&
\end{tikzcd}}

We seek the largest subgroup of \(A_g\) that normalizes \unsure{these|the} triple (72), and which \deleted{\unsure{defines}} \inserted{operates} therefore \deleted{\unsure{on}} the fibration \(\widetilde U_{g,\nu}\to U_{g,\nu}\) (where \(\widetilde U_{g,\nu}\) is the universal covering of \(U_{g,\nu}\) based at \(a_{g,\nu}\)); since the normalizers of \(A_{g,\nu}^{\circ}\) and \(A_{g,\nu+1}^{\circ}\) respectively are \(A_{g,\nu}\) and \(A_{g,\nu+1}\), the group in question must be contained in their intersection \(A_{g,\nu}^{\bullet}\), which normalizes \deleted{\unsure{also}} \(B_{g,\nu}=A_{g,\nu}^{\bullet}\cap A_{g,\nu}^{\circ}\). \unsure{It is} therefore this group that should be made to operate on the fibration (73) -- the subgroup \unsure{formed} \deleted{\illegible{7} \(U_{g,\nu}\)} \inserted{of the elements of \(A_{g,\nu}\)} whose operation

\scanpage{252}

on \(U_{g,\nu}\) is isotopic to the identity being \unsure{precisely} \(B_{g,\nu}\) -- so it is the group \unsure{\(\Gamma\)}
\[
 A_{g,\nu}^{\bullet}/B_{g,\nu}\simeq\Gamma_{g,\nu}
\]
that, as it should, operates up to isotopy. We would like to describe the group of all automorphisms of the topological fibration \(\widetilde U_{g,\nu}\) over \(U_{g,\nu}\) which will a priori be an extension of the group
\[
 \begin{aligned}
 A_{g,\nu}&=\operatorname{Aut}(U_{g,\nu})\quad\text{by}\quad
 \Pi_{g,\nu}=\pi_1(\unsure{U_{g,\nu}|\widetilde U_{g,\nu}},a_{g,\nu})\\
 &=\operatorname{Aut}_{U_{g,\nu}}(\widetilde U_{g,\nu}).
 \end{aligned}
\]

\deleted{\illegible{1}} This extension \inserted{is \uline{\unsure{split canonical\textsuperscript{ly}}}} \deleted{\unsure{split}} over \(A_{g,\nu}^{\bullet}\), \inserted{stabilizer at the base point \(a_{g,\nu}\)} (and we thus recover the operation of \(A_{g,\nu}^{\bullet}\) for the fibration (73). Now we already have an extension \(\Gamma'_{g,\nu+1}\) of \(\Gamma_{g,\nu}\) by \(\Pi_{g,\nu}\), whence by inverse image under \(A_{g,\nu}\to\Gamma_{g,\nu}=A_{g,\nu}/A_{g,\nu}^{\circ}\), an extension (which I am going to denote by \(\widetilde A_{g,\unsure{\nu}}\)) of \(A_{g,\nu}\) by \(\Pi_{g,\nu}\):
\[
 \widetilde A_{g,\nu}\simeq A_{g,\nu}\mathbin{\times}_{\Gamma_{g,\nu}}\Gamma'_{g,\nu+1}
 \tag{74}
\]
\(=\) quotient group of the subgroup \deleted{\unsure{\(H\)}} \(\widetilde A_{g,\nu}^{\unsure{\natural|\sharp}}\) \unsure{of} \(A_{g,\nu}\times A_{g,\nu}^{\bullet}\) formed of the pairs \((u,v)\), \(u\unsure{\equiv}v\) \deleted{\unsure{\(A_{g,\nu}^{\circ}\)}} modulo \(A_{g,\nu}^{\circ}\), by the subgroup \deleted{\unsure{\(B_{g,\nu}^{\circ}\)}} \(1\times A_{g,\nu+1}^{\circ}\).

We thus have two \inserted{(and = three!)} extension structures on \(\widetilde A_{g,\nu}\)

(75) \diagram{\begin{tikzcd}[ampersand replacement=\&][column sep=small]
1 \arrow[r] \& \Pi_{g,\nu} \arrow[r] \& \widetilde A_{g,\nu} \arrow[r,"p"] \& A_{g,\nu} \arrow[r] \& 1 \\
1 \arrow[r] \& A_{g,\nu}^{\circ} \arrow[r] \& \widetilde A_{g,\nu} \arrow[r] \& \Gamma'_{g,\nu+1} \arrow[r] \& 1
\end{tikzcd}}
\diagram{\begin{tikzcd}[ampersand replacement=\&][column sep=small]
1 \arrow[r] \& (A_{g,\nu}^{\circ}\unsure{\times|\rtimes}\Pi_{g,\nu}) \arrow[r] \& \widetilde A_{g,\nu} \arrow[r] \& \Gamma_{g,\nu} \arrow[r] \& 1
\end{tikzcd}}

\margin{\uline{N.B.} The second \unsure{is meant to express} that \(\widetilde A_{g,\nu}^{\circ}=A_{g,\nu}^{\circ}\) and \(\pi_0(\widetilde A_{g,\nu})\simeq\Gamma'_{g,\nu+1}\) (the canonical extension of \(\Gamma_{g,\nu}\) by \(\Pi_{g,\nu}\)).}
\deleted{\unsure{Now}}

\scanpage{253}
\mspage{127}

I would like to make \(\widetilde A_{g,\nu}\) operate on \(\widetilde U_{g,\nu}\), i.e. make \deleted{\illegible{1}} \(H=\widetilde A_{g,\nu}^{\unsure{\natural|\sharp}}\) operate \unsure{with operation} \deleted{in such a way that \((u,v)\)} trivial on \(A_{g,\nu+1}^{\circ}\), and this respecting the following conditions

\margin{\unsure{trivial datum} on \(1\times A_{g,\nu+1}^{\circ}\). Cf. generalization below in the margin.}

a) \deleted{\illegible{1}} (compatibility with \(\widetilde A_{g,\nu}\to A_{g,\nu}\)) the pair \((u,v)\in H\) operates \unsure{on} \(\widetilde U_{g,\nu}\) by an automorphism compatible with the automorphism \inserted{\(u\) of} \(U_{g,\nu}\) (we shall say that it is a \(u\)-automorphism).

b) (compatibility with \(\Pi_{g,\nu}\to\widetilde A_{g,\nu}\)) If \(u=1\), \inserted{[so \(v\in A_{g,\nu}^{\circ}\), so (since \(v\in A_{g,\nu}^{\bullet}\)) \(v\in B_{g,\nu}\)]} then the operation of \((u,v)=(1,v)\) on \(\widetilde U_{g,\nu}\) is none other than that defined by \(v\bmod B_{g,\nu}^{\circ}=A_{g,\nu+1}^{\circ}\), which is an element of \(\Pi_{g,\nu}\simeq\operatorname{Aut}(\widetilde U_{g,\nu}/U_{g,\nu})\), or again, it is that defined by \uline{translation} \inserted{\uuline{on the right}} by \(\unsure{v^{-1}}\).

\margin{specification of the operation on \(1\times B_{g,\nu}\).}

c) Compatibility with \deleted{\unsure{\(\widetilde A_{g,\nu}\to\Gamma'_{g,\nu+1}\)}} the operation already \deleted{\illegible{1}} \inserted{obtained above} of \(A_{g,\nu}^{\bullet}\) on \(\widetilde U_{g,\nu}\): if \((u,v)\in H\) is such that \(u=v\) (so \(u=v\in A_{g,\nu}^{\bullet}\)) then \(\unsure{(u,v)}=(u,u)\) operates \unsure{via} the inner automorphism defined by \(u\).

\margin{specification of the operation on \(\delta(A_{g,\nu}^{\bullet})\) (\(\delta\): diagonal arrow).}

d) \inserted{Compatibility with \(A_{g,\nu}^{\circ}\to\widetilde A_{g,\nu}\):} The operation of \(A_{g,\nu}^{\circ}\) on \(\widetilde U_{g,\nu}\) is continuous (which, together with a), determines this operation uniquely \deleted{\illegible{2}} \inserted{and} the existence \inserted{\unsure{a priori of such an operation}} (follows from \(\pi_1(A_{g,\nu}^{\circ})=0\))).

\margin{specification of the operation on \((A_{g,\nu}^{\circ})\times 1\).}

\unsure{every} element \inserted{\((u,v)\)} of \(H\) can be written uniquely as a product of an element \inserted{\((v,v)\)} of \(\delta(A_{g,\nu}^{\bullet})\), with an element \inserted{\((v^{-1}u,1)\)} of \(A_{g,\nu}^{\circ}\times 1\) (the first group normalizing the second)

\margin{\uline{N.B.} The operation of \(A_{g,\nu}^{\circ}\times 1\) is \unsure{described} most simply by \uline{translations} of \unsure{\(A_{g,\nu}^{\circ}\)} on the \unsure{homogeneous} space \(\widetilde U_{g,\nu}\) of \(A_{g,\nu}^{\circ}\).}

\scanpage{254}

and an operation of this \(\frac12\) direct product is given \inserted{indeed} by \deleted{\illegible{1}} specifying operations of the two factor groups, satisfying a compatibility condition, which is verified here by transport of structure. We have indeed defined an operation of \(H=\widetilde A_{g,\nu}^{\unsure{\natural|\sharp}}\) on the fibration (73), \deleted{\illegible{2}} satisfying conditions c) and d) by the very construction, it remains to verify a) and b).

Now for a), it suffices to verify it separately for the \unsure{elements} of \(H\) in \(\delta(A_{g,\nu}^{\bullet})\) and in \(A_{g,\nu}^{\circ}\times 1\), where this is trivial by construction in both cases. It remains to verify b), i.e. to make explicit the operation of an \unsure{element} \((1,v)\), \(v\in B_{g,\nu}\), which we write as
\[
 (\qquad (1,v)=\deleted{(\unsure{v^{-1}},1)}(v,v)(v^{-1},1)
\]
hence
\[
 \rho\deleted{_{\unsure{\widetilde A}_{g,\nu}}}(1,v)=\rho(v,v)\rho(v^{-1},1).
\]

Now \(\rho(v,v)\) is induced by \(\operatorname{int}(v)\), \inserted{i.e. \(x\mapsto vxv^{-1}\)} and \(\rho(v^{-1},1)\) by left translation: \(x\mapsto v^{-1}x\), so the composite operates by \(x\mapsto xv^{-1}\). Phew!

\scanpage{255}
\mspage{128}

We now assume that \(X_g\) is equipped with a \(C^\infty\) structure, and we replace in the considerations \deleted{Let} above the groups of homeomorphisms by groups of diffeomorphisms.

\margin{\uline{N.B.} A \(C^{\unsure{\nu|r}}\) structure with \(\unsure{\nu|r}\geq 2\) would suffice}
\[
 \begin{aligned}
 E_g={}&\text{set of \deleted{complex} \inserted{conformal} structures}\\
 &\text{on }X_g\text{ compatible with its}\\
 &\text{structure }C^\infty
 \end{aligned}
 \tag{76}
\]

\deleted{\unsure{Then}} We see at once that \(E_g\) is an \uline{\(\infty\)-connected} topological space, as a quotient of the \(\infty\)-connected space \inserted{(a convex set)} of Riemannian structures, by the \deleted{\unsure{group}} \(\infty\)-connected group of \(C^\infty\) maps from \(X_g\) to \(\mathbb R^{+*}\). On \(E_g\) the group \(A_g\) operates -- but of course \(E_g\) is no longer a homogeneous space. Let us note at once
\[
 \begin{aligned}
 E_g/A_g\simeq{}&\text{set of isomorphism classes}\\
 &\text{of \deleted{algebraic curves} \inserted{surfaces} \deleted{\unsure{compact}}}\\
 &\text{conformal compact orientable}\\
 &\text{of genus }g
 \end{aligned}
 \tag{77}
\]

And if we choose one of the two orientations of \(X_g\), so that \(E_g\) is identified with the set of \uline{complex} structures on \(X_g\) compatible with its \(C^\infty\) structure \inserted{and its orientation}, then \(E_g/A_g\) is identified with the space of isomorphism classes of compact complex curves \inserted{(nonsingular) connected} of genus \(g\), \uline{modulo} passage to the complex conjugate. On the other hand

\scanpage{256}

\[
 \begin{aligned}
 E_g/A_g^+\simeq{}&\text{set of isomorphism classes}\\
 &\text{of curves }\mathbb C\text{-algebraic \inserted{smooth} projective}\\
 &\text{connected of genus }g.
 \end{aligned}
 \tag{78}
\]

More generally

\[
 \begin{aligned}
 E_g/A_{g,\nu}\simeq{}&\text{set of isomorphism classes}\\
 &\text{of \deleted{structures} \deleted{on the} compact surfaces}\\
 &\text{\inserted{conformal} equipped \unsure{with multipunctures|with multipuncturings}}\\
 &\text{of type }(g,\nu)
 \end{aligned}
 \tag{79}
\]

\[
 \begin{aligned}
 E_g/A_{g,\nu}^{+}\simeq{}&\text{set of isomorphism classes}\\
 &\text{of curves \deleted{\unsure{complex}} \inserted{\(\mathbb C\) algebraic}}\\
 &\text{[smooth projective connected}\\
 &\text{of genus }g\text{] \uline{equipped} with a}\\
 &\text{multipuncturing, of type}\\
 &(g,\nu)
 \end{aligned}
 \tag{80}
\]
\margin{if \(X_g\) oriented.}

These are the ``coarse moduli spaces'' (Mumford's ``coarse moduli''), which we may denote by \(M_g^{\unsure{\natural|\sharp}}\) and \(M_{g,\nu}^{\unsure{\natural|\sharp}}\), and which are indeed too coarse for the most important geometric uses. \deleted{But} Much more interesting are the quotients \(E_g/A_{g,\nu}^{\circ}\), \illegible{1}

\scanpage{257}
\mspage{129}

\[
 \widetilde M_{g,\nu}=E_g/A_{g,\nu}^{\circ}
 \qquad
 (\widetilde M_g=\widetilde M_{g,0}=E_g/A_g^{\circ})
 \tag{81}
\]
on which the group \deleted{\illegible{1}}
\[
 \Gamma_{g,\nu}=A_{g,\nu}/A_{g,\nu}^{\circ}
\]
operates by passage to the quotient of the operation of \(A_{g,\nu}\).

\deleted{\illegible{1}} The space \(\widetilde M_{g,\nu}\) \inserted{(with the operation of \(\Gamma_{g,\nu}\))} is called \uline{the Teichmüller space} of type \(g,\nu\). Of course we recover \(M_{g,\nu}^{\unsure{\natural|\sharp}}\) from \(\widetilde M_{g,\nu}\) and \margin{\deleted{(82)}} the operation of \(\Gamma_{g,\nu}\) on it, by
\[
 M_{g,\nu}^{\unsure{\natural|\sharp}}=\widetilde M_{g,\nu}/\Gamma_{g,\nu}
 \tag{\unsure{82}}
\]

\uline{Theorem} (Teichmüller) The Teichmüller space \(\widetilde M_{g,\nu}\) is homeomorphic to \deleted{\(\mathbb R\)} \(\mathbb C^\mu\), where \(\mu=3g-3+\nu\) (\inserted{in the} anabelian cases \(2g+\nu\geq 3\)), and \(\mu=3g-3+(\unsure{\nu}+\delta)\)
\[
 \delta=\dim_{\mathbb C}G
\]
in the general case, \(G\) being \unsure{group} of algebraic automorphisms of a complex \(U_{g,\nu}\) (so \(\delta=1\) in the \unsure{``limiting''} cases \((1,0)\) and \((0,2)\), and \inserted{more \unsure{generally}} \(\delta\) increases by 1 each time \unsure{for} \(g\) fixed, and \((g,\nu-1)\) abelian, one passes from \(\nu\) to \(\nu-1\)), hence

\margin{We can also say that \(\mu=3g-3+\nu'\), where \(\nu'=\nu\) in the \unsure{anabelian} case, and \unsure{where \(\nu'\)} in the ``general'' case is the least integer \(\geq\nu\) such that \((g,\nu')\) is \unsure{``anabelian''}.}

\scanpage{258}

\[
 \left\{
 \begin{aligned}
 \mu&=1\quad\text{in the case }(1,0)\\
 \mu&=0\quad\text{in the cases }(0,2),(0,1),(0,0),
 \end{aligned}
 \right.
 \tag{\unsure{\(82_1\)|\(82'\)}}
\]
i.e. \(\widetilde M_{g,\nu}\) is in these cases reduced to a point, i.e. the action of \(A_{g,\nu}^{\circ}\) on \(E_g=E_0\) is \deleted{already} transitive -- (and these are \inserted{with the anabelian case \((0,3)\)} the only three cases where this is so, \deleted{\illegible{4}}).

\margin{i.e. on surfaces \unsure{oriented} \unsure{spherical with} \(\unsure{\nu\leq 3}\), all complex structures are \uline{isotopic}.}

Starting from \(\widetilde M_{0,3}\), or from \(\widetilde M_{1,1}\) or from \(\widetilde M_{g,0}=\widetilde M_g\) with \(g\geq 1\), when \(\nu\) increases the ``complex dimensions'' of \(\widetilde M_{g,\nu}\) increase by the same amount (On the other hand \(\widetilde M_{1,0}\unsure{\longleftarrow}\widetilde M_{1,1}\) is a homeomorphism \deleted{\unsure{because}\ \(A_{1,0}^{\circ}\unsure{\simeq}A_{1,1}^{\circ}\) \illegible{1}}).

We can make Teichmüller's theorem more precise as follows

\uline{Corollary} In the \uline{anabelian} case, \(A_{g,\nu}^{\circ}\) operates freely on \(E_g\), so that \(E_g\) becomes a torsor over \(\widetilde M_{g,\nu}\), with structure group \(A_{g,\nu}^{\circ}\).

We deduce from this that \(\widetilde M_{g,\nu}\) plays the role of a \uline{classifying space} for \(A_{g,\nu}^{\circ}\), and that
\[
 \pi_{\unsure{i|\nu}}(\widetilde M_{g,\nu})
 \xrightarrow{\sim} \pi_{\unsure{\nu-1|i-1}}(A_{g,\nu}^{\circ})
 \tag{83}
\]
and the fact that \(\widetilde M_{g,\nu}\) is \(\infty\)-connected (contained in Teichmüller's theorem) is equivalent to

\scanpage{259}
\mspage{130}

the fact that \(A_{g,\nu}^{\circ}\) is so -- which is the ``well-known theorem'' recalled in number 19.

I do not dwell here on the interpretation of the points of \(\widetilde M_{g,\nu}\) as isomorphism classes of complex curves, equipped \deleted{with} \inserted{with a} suitable ``Teichmüller rigidification'', \unsure{and the rigidified moduli spaces point of view}, which makes it possible to verify a priori that \(\widetilde M_{g,\nu}\) is equipped with the structure of a nonsingular complex manifold -- but already the fact that \(\widetilde M_{g,\nu}\) is \uline{simply connected} (which can \unsure{express} \unsure{by} interpreting \(\widetilde M_{g,\nu}\) as the universal covering of the \uline{modular topos} \(M_{g,\nu}\)) is quite deep, which does not seem to be able to fit into the framework of topology (or of differential topology \((C^\infty)\) alone...

\margin{\uline{N.B.} To prove that the \(\unsure{\widetilde M_{g,\nu}}\) are \(\infty\)-connected, \unsure{we are easily reduced} to the case of \(M_{g,0}\) \unsure{(if \(g\geq 2\))}, or of \(M_{1,1}\) (punctured elliptic case), by using \unsure{the fibrations} \(A_{g,\nu}^{\circ}/A_{g,\nu+1}^{\circ}\simeq\widetilde U_{g,\nu}\) (which is \(\infty\)-connected...). The case \(g=1\) is \unsure{moreover easy and} well understood...}

\scanpage{260}

Let us note that the subgroup inclusions
\[
 A_{g,\nu+1}^{\circ}\subset A_{g,\nu}^{\circ}\cdots\subset A_{g,0}^{\circ}=A_g^{\circ}
\]
define a tower of continuous maps
\[
 \cdots\longrightarrow\widetilde M_{g,\nu+1}\longrightarrow\widetilde M_{g,\nu}
 \longrightarrow\cdots\longrightarrow\widetilde M_{g,0}=\widetilde M_g
 \tag{84}
\]

where
\[
 \widetilde M_{g,\nu+1}\longrightarrow\widetilde M_{g,\nu}
 \tag{85}
\]
is \inserted{\uline{for \((g,\nu)\) anabelian}} a fibration with fiber \(A_{g,\nu}^{\circ}/A_{g,\nu+1}^{\circ}\simeq\widetilde U_{g,\nu}\).

[it therefore has contractible fiber, and the \(\infty\)-connectedness of \(\widetilde M_{g,\nu+1}\) is therefore equivalent to that of \(\widetilde M_{g,\nu}\) -- which brings us back to the case of \(\widetilde M_g\) if \(g\geq 2\), of \(M_{\text{\deleted{\(g\)}}1,1}\) and of \(M_{0,3}\) if \(g\unsure{=|\neq|\geq}1\) or 0].

In the abelian case \((g,\nu)\), we find
\[
 \widetilde M_{g,\nu+1}\simeq\widetilde M_{g,\nu}\qquad
 (\text{if }(g,\nu)\text{ abelian})
 \tag{85}
\]
which \deleted{\unsure{follows precisely|stems, follows} from the fact already} \inserted{means} \deleted{noted} that in these cases, if we have two complex structures \inserted{\(\alpha,\beta\)} on \(X\) which are congruent by \inserted{\(u\in\)} \(A_{g,\nu}^{\circ}\), \inserted{(i.e. \(\beta=u.\alpha\))} they are = congruent by \(A_{g,\nu+1}^{\circ}\).

Indeed, if \(G\) is \unsure{the} \inserted{\unsure{identity|neutral} component of the} group of complex automorphisms of \(U_{g,\nu}\) for \(\beta\), we can write \(u=gv\) with \(g\in G\), \(v\in A_{g,\nu+1}^{\circ}\) hence \(\beta=gv.\alpha\) hence (since \(g^{-1}\beta=\beta\)) \(\beta=v\alpha\), q.e.d.
